\section{Results}
\label{sec:results}

\subsection{Shift falls into a few regimes, and the common one is a threshold shift}
\label{sec:results-regimes}

Figure~\ref{fig:regime-map} places each of the 21 shifted source--target pairs by the shift of its predicted-label
marginal, which needs no target labels, and by its AUROC retention; Table~\ref{tab:regime-map-source-transfer}
(Appendix~\ref{sec:appendix-results}) lists every coordinate and the best correction per pair. The post-hoc thresholds
are: retention $\ge 0.95$, $|\log T^*_Q/T_P| \le 0.15$ and marginal shift $\le 0.10$ for no correction; a ratio of
at least 0.25 with a marginal shift $\le 0.10$ for a temperature shift; and a marginal shift $\ge 0.10$ with
target TS better than a constant in most runs and a finite oracle temperature for a threshold shift.

\begin{figure}[t]
\centering
\includegraphics[width=\linewidth]{regime_map_source_transfer.pdf}
\caption{Regime map of the 21 shifted source--target pairs. Each point is a target of one source task (mean over the
converged cross-entropy runs of that source, each after its own source temperature scaling). Horizontal: total-variation
shift of the predicted-label marginal, computed without target labels. Vertical: target AUROC divided by in-domain
AUROC. Hollow: target temperature scaling is on average no better than a constant predictor. Cross: unbounded oracle
temperature in every run.}
\label{fig:regime-map}
\end{figure}

\paragraph{Mild varieties need nothing; pure temperature shifts are rare.}
Encoders trained on BD-SHS keep their ranking on the Barishal and Noakhali dialects (retention 0.994 and 0.990, marginal
shift at most 0.035), and SentNoB encoders keep 0.957 of their ranking on Sylheti; no correction improves on source TS
there by more than 0.0035 nats. Only two pairs show a temperature shift with stable predicted proportions: the
Chittagong dialect ($\log T^*_Q/T_P = 0.297$, marginal shift 0.029) and BD-SHS as a target of the OLID-trained models
(0.407 and 0.082).

\paragraph{The common failure keeps the ranking and moves the threshold.}
Twelve of the 21 pairs, from all four source tasks, keep a usable ranking (retention 0.70--0.98; target TS beats the
constant predictor in 88--100\% of runs) while their predicted marginal moves by 0.10--0.46. They are natural Banglish,
Bangla-script code-mixing, MixTox-SA and English OLID for the BD-SHS models; OffMix-3L, both versions of Code-mixed Chaos
and MixTox-SA for the OLID models; SentMix-3L and Hinglish for the SentNoB models; and SentMix-3L and SentNoB for the
Hinglish models. The movement is not a change of the class prior: on Banglish, Bangla-script Chaos and SentMix-3L the
true class proportions differ from the source by a total variation of 0.007--0.021, while the predicted ones move by
0.227--0.404. In every one of the 12 pairs a bias correction raises macro-F1 in 98--100\% of runs
(Section~\ref{sec:results-threshold}): the decision threshold is in the wrong place.

\paragraph{Weak ranking and collapse.}
Spanglish, a language absent from both sentiment sources, keeps AUROC 0.609 (SentNoB models) and 0.648 (Hinglish
models), but target TS is no better than a constant in 56\% and 73\% of runs; Sylheti behaves the same way for the
Hinglish models (53\%). A discrimination collapse occurs for exactly one pair, BD-SHS models on OffMix-3L: AUROC 0.496
(0.387--0.589 across runs) with an unbounded oracle temperature in every run.

\paragraph{Against the pre-registered reading.}
We had decided to call the findings independent of the source language only if the English and Hinglish sources showed
the same regimes and the same ranking of corrections as the Bangla ones. The threshold shift, with the same ranking of
corrections, appears for all four sources and the temperature shift for a Bangla and the English source. No-correction
appears only for Bangla sources, the only ones with a same-language mild variety among their targets, and collapse
appears for one Bangla source and disappears when the source uses the target's label definition
(Section~\ref{sec:results-causes}). The strong form of the prediction is not met, and we restrict these two regimes to
the Bangla sources.

\subsection{Under a threshold shift, temperature is the wrong knob}
\label{sec:results-threshold}

Table~\ref{tab:correction-ablation-by-regime} adds the two components of a correction one at a time. In the
threshold-shift regime, shrinkage TS with 50 target labels lowers NLL from 0.892 to 0.728 but leaves macro-F1 at
0.527, as it must. Label-free batch calibration, which shifts only the class bias, raises macro-F1 to 0.620; MAP BCTS
with 50 labels, which adjusts both, reaches 0.594 and the lowest NLL (0.684). The temperature is not merely
insufficient but the wrong axis: label-free batch calibration has lower NLL than the oracle temperature fitted on all
target labels, for example 0.848 against 1.008 for the Hinglish models on SentMix-3L and 0.573 against 0.647 for the
OLID models on OffMix-3L. Training-time remedies that act like a temperature behave the same way: focal loss and label
smoothing lower the uncalibrated ECE and NLL on most shifted conditions, but after source TS neither is significantly
better than cross-entropy on any condition, label smoothing is significantly worse on four of the five shifted
offensive conditions, and no objective, scale or architecture lifts the BD-SHS models' AUROC on OffMix-3L above 0.546
(Appendix~\ref{sec:appendix-objectives}).

\begin{table}[t]
\centering
\small
\caption{Which part of the correction matters in each regime (all converged cross-entropy runs of the four source tasks; means over runs and cells of the regime). Shrinkage TS only rescales the temperature and cannot change the decision; batch calibration only shifts the class bias and uses no labels; MAP BCTS does both with 50 target labels. Last column: share of run $\times$ cell units whose NLL is below source TS. Cells: number of (source task, target) pairs in the regime, as assigned in the regime map.}
\label{tab:correction-ablation-by-regime}
\resizebox{\textwidth}{!}{%
\begin{tabular}{lrlllrrr}
\toprule
Regime & Cells & Correction & Adjusts & Labels & Macro-F1 & NLL & NLL $<$ source TS \\
\midrule
R0 no correction & 3 & Source TS & none & no & 0.803 & 0.442 & -- \\
 &  & Shrinkage TS & temperature & 50 & 0.803 & 0.450 & 0.00 \\
 &  & Batch calibration & bias & no & 0.802 & 0.455 & 0.63 \\
 &  & MAP BCTS & temperature + bias & 50 & 0.802 & 0.449 & 0.17 \\
\midrule
R1 temperature shift & 2 & Source TS & none & no & 0.792 & 0.464 & -- \\
 &  & Shrinkage TS & temperature & 50 & 0.792 & 0.453 & 0.72 \\
 &  & Batch calibration & bias & no & 0.797 & 0.454 & 0.94 \\
 &  & MAP BCTS & temperature + bias & 50 & 0.797 & 0.452 & 0.61 \\
\midrule
R2 threshold shift & 12 & Source TS & none & no & 0.527 & 0.892 & -- \\
 &  & Shrinkage TS & temperature & 50 & 0.527 & 0.728 & 0.90 \\
 &  & Batch calibration & bias & no & 0.620 & 0.699 & 0.98 \\
 &  & MAP BCTS & temperature + bias & 50 & 0.594 & 0.684 & 0.94 \\
\midrule
R3 weak ranking & 3 & Source TS & none & no & 0.376 & 1.118 & -- \\
 &  & Shrinkage TS & temperature & 50 & 0.376 & 1.036 & 1.00 \\
 &  & Batch calibration & bias & no & 0.429 & 1.077 & 0.59 \\
 &  & MAP BCTS & temperature + bias & 50 & 0.391 & 1.014 & 1.00 \\
\midrule
R4 collapse & 1 & Source TS & none & no & 0.367 & 1.299 & -- \\
 &  & Shrinkage TS & temperature & 50 & 0.367 & 0.715 & 1.00 \\
 &  & Batch calibration & bias & no & 0.473 & 0.740 & 1.00 \\
 &  & MAP BCTS & temperature + bias & 50 & 0.448 & 0.704 & 1.00 \\
\bottomrule
\end{tabular}
}
\end{table}


\paragraph{When label-free bias correction backfires.}
Batch calibration equalises the predicted class proportions, so it assumes a balanced target. To test this assumption we
resampled every target of the two offensive source tasks to a fixed positive share $\pi_+$
(Table~\ref{tab:prior-shift-ablation}). At $\pi_+ = 0.5$ batch calibration comes close to MAP BCTS (macro-F1 0.699
against 0.709). At $\pi_+ = 0.1$, the typical situation in content moderation, it is harmful: macro-F1 drops to 0.631
from 0.654 with source TS, NLL rises from 0.31 to 0.56, and it lowers NLL in none of the run--target units. A class prior
estimated from 50 labels repairs it (0.661 and 0.27) and performs like the true prior; MAP BCTS does as well (0.665 and
0.25). At $\pi_+ = 0.9$ the prior-aware variant reaches 0.615 against 0.482 for the uniform one.

\begin{table}[t]
\centering
\small
\caption{Batch calibration under class-prior shift (offensive source tasks, BD-SHS and OLID; every target resampled to a positive share $\pi_+$; means over runs, targets and 5 resamples; macro-F1 / NLL). Uniform batch calibration assumes $\pi_+ = 0.5$; a prior estimated from 50 labels removes that assumption and is close to the true-prior reference.}
\label{tab:prior-shift-ablation}
\begin{tabular}{lccccc}
\toprule
Method & $\pi_+ = 0.1$ & $\pi_+ = 0.3$ & $\pi_+ = 0.5$ & $\pi_+ = 0.7$ & $\pi_+ = 0.9$ \\
\midrule
Source TS & 0.654 / 0.31 & 0.666 / 0.56 & 0.630 / 0.81 & 0.561 / 1.06 & 0.434 / 1.30 \\
Batch calibration (uniform prior) & 0.631 / 0.56 & 0.716 / 0.50 & 0.699 / 0.54 & 0.631 / 0.65 & 0.482 / 0.83 \\
Batch calibration + prior from 50 labels & 0.661 / 0.27 & 0.691 / 0.47 & 0.698 / 0.54 & 0.674 / 0.51 & 0.615 / 0.32 \\
Batch calibration + true prior (reference) & 0.661 / 0.26 & 0.692 / 0.46 & 0.699 / 0.54 & 0.682 / 0.51 & 0.612 / 0.31 \\
MAP BCTS (50 labels) & 0.665 / 0.25 & 0.686 / 0.45 & 0.709 / 0.51 & 0.636 / 0.46 & 0.572 / 0.26 \\
\bottomrule
\end{tabular}
\end{table}


\paragraph{The predicted marginal cannot tell when to correct.}
A threshold on the predicted-marginal shift, fitted on the Bangla sources and applied to the OLID and Hinglish sources,
identifies the pairs that need a bias correction with AUROC 0.845 and balanced accuracy 0.796. Yet 92\% of the held-out
units need one, so always correcting after a detected shift is more accurate (0.919 against 0.778). The same statistic
does not identify the units in which even target TS is no better than a constant (balanced accuracy 0.524 on the
Hinglish source). Recognising weak ranking and collapse requires labels, as Proposition~\ref{prop:impossibility}
anticipates.

\subsection{What causes the collapse: script in part, label definition above all}
\label{sec:results-causes}

\paragraph{Script.}
Code-mixed Chaos exists in Banglish and in a row-aligned Bangla-script transliteration with identical labels. Changing
only the script raises encoder AUROC by 0.212--0.225 for IndicBERT (four training variants), 0.159 for mDeBERTa-v3,
0.071--0.096 for XLM-R-base, 0.043 for MuRIL and 0.022--0.036 for BanglaBERT, while the LLMs change by $-0.001$ to $+0.055$
(Figure~\ref{fig:script}). Script coverage explains part of the encoders' difficulty with romanised text, but not the
OffMix-3L collapse: MuRIL, which was pretrained with transliterated text, still ranks OffMix-3L at 0.546.

\paragraph{Label definition.}
OffMix-3L derives its notion of offensiveness from OLID \citep{goswami2023offmix}, which counts any
non-acceptable language as offensive, including untargeted profanity \citep{zampieri2019olid}; BD-SHS annotates hate
speech \citep{romim2022bdshs}. Training the same three architectures on
English OLID instead of BD-SHS raises OffMix-3L AUROC from 0.519 to 0.819 (pooled difference $+0.300$, 95\% CI
$[0.275, 0.326]$) while lowering it on Banglish Chaos; the pre-registered difference-in-differences is $+0.378$
($[0.351, 0.404]$, $p_{\mathrm{Holm}} < 0.001$; Table~\ref{tab:cross-definition-contrasts},
Figure~\ref{fig:cross-definition}). The crossover is complete: BD-SHS-trained models rank BD-SHS test at 0.972 and OLID
test at 0.686, OLID-trained models at 0.779 and 0.869. The collapse is therefore a source--target mismatch of the label
concept, shown by intervention and without relying on our own judgement of the labels. Matching the definition restores
the ranking but not the threshold: OLID-trained models label 8.4\% of OffMix-3L as offensive, against 47.8\% in the gold
labels.

\begin{figure}[t]
\centering
\begin{minipage}[t]{0.46\linewidth}
\centering
\includegraphics[width=\linewidth]{script_intervention_auroc.pdf}
\caption{Same comments, two scripts: AUROC on Code-mixed Chaos in Banglish and in Bangla script.}
\label{fig:script}
\end{minipage}\hfill
\begin{minipage}[t]{0.49\linewidth}
\centering
\includegraphics[width=\linewidth]{cross_definition_auroc.pdf}
\caption{Same architectures, two label definitions: encoders trained on BD-SHS or on English OLID.}
\label{fig:cross-definition}
\end{minipage}
\end{figure}

\begin{table}[t]
\centering
\small
\setlength{\tabcolsep}{4pt}
\caption{Cross-definition test: AUROC of encoders trained on BD-SHS vs on English OLID (XLM-R-base, MuRIL, mDeBERTa-v3; 5 seeds each). Differences are OLID-trained minus BD-SHS-trained; CIs from a paired item bootstrap; $p$ from a stratified permutation test, Holm-adjusted over the seven contrasts. The last row is the pre-registered difference-in-differences.}
\label{tab:cross-definition-contrasts}
\begin{tabular}{lrrrcr}
\toprule
Condition & BD-SHS-trained & OLID-trained & Difference & 95\% CI & $p_{\mathrm{Holm}}$ \\
\midrule
BD-SHS test & 0.972 & 0.779 & -0.193 & [-0.203, -0.183] & $<$0.001 \\
OLID test & 0.686 & 0.869 & +0.183 & [+0.152, +0.214] & $<$0.001 \\
OffMix-3L & 0.519 & 0.819 & +0.300 & [+0.275, +0.326] & $<$0.001 \\
Chaos (Banglish) & 0.783 & 0.706 & -0.077 & [-0.084, -0.071] & $<$0.001 \\
Chaos (Bangla script) & 0.879 & 0.854 & -0.025 & [-0.031, -0.019] & $<$0.001 \\
MixTox-SA & 0.712 & 0.744 & +0.033 & [+0.024, +0.041] & $<$0.001 \\
OffMix-3L minus Chaos (Banglish) & -0.264 & 0.113 & +0.378 & [+0.351, +0.404] & $<$0.001 \\
\bottomrule
\end{tabular}
\end{table}


\paragraph{OffMix-3L is also artifact-laden.}
A character n-gram classifier trained on 50 OffMix-3L items reaches AUROC 0.976, and label-free writing-style clusters
explain much of its label (adjusted mutual information 0.47--0.81 for 5 to 50 clusters, against 0.19--0.28 on BD-SHS).
The comparisons above survive this artifact: within five clusters the zero-shot LLMs still rank offensive items above
non-offensive ones (AUROC 0.910--0.950) whereas the BD-SHS encoders do not (0.470 on average), and the OLID-trained
encoders do (Appendix~\ref{sec:appendix-artifacts}). Because the natural code-mixed targets carry much weaker artifacts,
we treat them, not OffMix-3L, as the main evidence about code-mixing.

\subsection{Detecting the shift and choosing a system}
\label{sec:results-systems}

\paragraph{Confidence does not reveal the shift; a domain classifier does.}
A character n-gram domain classifier separates in-domain from target items with AUROC 0.982--1.000 on every original
condition. The models' own uncertainty does not: the encoders reach 0.504--0.686 on average, and the LLMs' uncertainty is
inverted on tri-lingual code-mixing (0.352 offensive, 0.318 sentiment), meaning that they are more confident on the
shifted inputs. Under collapse, abstaining on low confidence gives no protection: on OffMix-3L the encoders' selective
risk at 80\% coverage is 0.476, against an overall error of 0.475.

\paragraph{LLMs keep the ranking that encoders lose; encoders win in domain.}
On identical items (Figure~\ref{fig:encoder-llm}), zero-shot LLMs rank OffMix-3L at 0.857--0.941 against 0.481 for the
BD-SHS encoders, Banglish at 0.886--0.970 against 0.773, MixTox-SA at 0.795--0.831 against 0.687, and Hinglish at
0.843--0.890 against 0.735. In domain and on the dialects the encoders are better (BD-SHS test 0.970 against
0.758--0.854). A guided-completion test found no near-verbatim completion of any test set, and Llama-3.1-8B, which shows
no completion signal at all, ranks OffMix-3L as well as the other LLMs (Appendix~\ref{sec:appendix-contamination}).

\begin{figure}[t]
\centering
\includegraphics[width=\linewidth]{encoder_llm_auroc_by_condition.pdf}
\caption{Encoders and zero-shot LLMs on identical items, after each system's own source temperature scaling.}
\label{fig:encoder-llm}
\end{figure}

\paragraph{How to spend 50 labels.}
The LLMs' rankings come with a shifted threshold. On OffMix-3L, MAP BCTS raises zero-shot LLM macro-F1 from 0.470 to
0.752 with 25 labels and 0.807 with 50, and conformal risk control at an FNR of 0.1 then costs an FPR of 0.32--0.36 for
three of the four LLMs (0.66 for Qwen3-8B), against 0.94 for the encoders. Spending the same 50 labels as in-context
demonstrations improves the ranking (AUROC 0.966 against 0.913) but gives lower macro-F1 (0.693 against 0.794) and NLL
2.005 against 0.429. Fine-tuning the encoder on the target labels cures the collapse only with about 200 labels (AUROC
0.505, 0.621, 0.743 and 0.894 with 25, 50, 100 and 200 labels; Figure~\ref{fig:few-label}). With up to 100 labels the
better use of them is to route to a system whose ranking survives and to correct its bias.

\begin{figure}[t]
\centering
\includegraphics[width=\linewidth]{few_label_adaptation_curves.pdf}
\caption{Ways of spending $n$ target labels, macro-F1 on the target test items. Dashed line: encoder with source TS.
The character n-gram classifier is trained on the same labels and shows how much of each label is predictable from
surface features.}
\label{fig:few-label}
\end{figure}
